\def\Ponderation{33,33}
\def\SigleNumero{STT-1920}
\def\NomCours{Méthodes statistiques}
\def\DateEvaluation{25 mars 2026}
\def\HeureEvaluation{13 h 30 à 15 h 20}
\def\Session{}
\def\NomEvaluation{Examen 2}

% Ne rien mettre entre les accolades si l'enseignant (ou l'enseignante) est aussi le coordonnateur (ou la coordonnatrice)
\def\Coordonnatrice{}
\def\Coordonnateur{}

% Ne rien mettre entre les accolades s'il y a plus d'une section
\def\Enseignant{Julien Miron}
\def\Enseignante{}

% Ne rien mettre entre les accolades s'il s'agit d'un cours à section unique
% Mettre le nom de chacun des enseignants et enseignantes entre les accolades
% Une case à cocher sera générée pour chacune des sections
\def\SectionA{}
\def\SectionB{}
\def\SectionC{}
\def\SectionD{}


\def\Directives{
\begin{enumerate}
\item Matériel autorisé:
\begin{itemize}
\item Calculatrice scientifique {\bf autorisée par la faculté des sciences et de génie};
\item Votre aide-mémoire d'une feuille recto-verso, écrit à la main.
\end{itemize}
\item Assurez-vous que cette évaluation comporte \numquestions ~exercices répartis sur \numpages{}~pages, pour un total de \numpoints{}~points.\
\item \textbf{Veuillez ne pas dégrafer votre examen.}
\item Vous devez\textbf{ justifier toutes vos réponses} par une démarche claire, sauf pour les questions à choix de réponse.\
\item Pour les questions à choix de réponse et les vrais ou faux, vous pouvez mettre un X, un crochet ($\checkmark$) ou remplir la case.\
\item Vous devez laisser tout ce qui ne servira pas durant l'examen à l'avant de la salle AVANT de rejoindre votre place.\
\item Vous n'avez droit à AUCUN appareil électronique en votre possession durant l'examen.
\item Une fois assis à votre place, vous devez demandez l'autorisation pour vous lever, même si l'examen n'est pas encore commencé.
\item Dans les 5 dernières minutes de l'examen, il sera interdit de se lever avant que TOUTES les copies n'aient été ramassées.
\end{enumerate}
\begin{itemize}
\item[$\rightarrow$] Le non-respect des consignes 7-8-9 entraînera une pénalité de 10\%.
\end{itemize}
}

\newcommand{\affichepoints}{}     % ← AVEC points
% \renewcommand{\affichepoints}{on} % ← SANS points (non vide)

\documentclass{Evaluation}
\usepackage{multicol,graphicx}

\noprintanswers



\begin{document}

\pageblanche

%%%%%%%%%%%%%%%% QUESTION %%%%%%%%%%%%%%%%
\begin{questions}

\question[3]\label{interp}
La durée de vie (en heures) d'un certain type d'ampoule est modélisée par une loi normale de moyenne $\mu$ inconnue et de variance $\sigma^2$ connue. À partir d'un échantillon de taille $n=16$, on calcule un intervalle de confiance à 95\% pour $\mu$ et on obtient $[980;\, 1020]$. Quel énoncé est une interprétation \textbf{correcte} de cet intervalle?

\begin{checkboxes}
    \choice 95\% des ampoules de ce type ont une durée de vie entre 980 et 1020 heures.
    \choice La probabilité que $\mu$ soit entre 980 et 1020 heures est de 0,95.
    \correctchoice On a confiance à 95\% que l'intervalle $[980;\, 1020]$ contient la vraie valeur de $\mu$.
    \choice Si on fabrique 100 ampoules, 95 d'entre elles dureront entre 980 et 1020 heures.
\end{checkboxes}

\begin{solution}
Un IC à 95\% signifie que, si on répétait l'expérience un grand nombre de fois, 95\% des intervalles calculés contiendraient $\mu$. On a donc confiance à 95\% que cet intervalle particulier contient $\mu$. Les autres choix confondent l'intervalle pour $\mu$ avec la distribution individuelle des observations ou attribuent une probabilité au paramètre fixe $\mu$.
\end{solution}


\question[3] Laquelle des affirmations suivantes concernant la marge d'erreur d'un intervalle de confiance pour la moyenne $\mu$ est \textbf{vraie}?

\begin{checkboxes}
    \choice Si on augmente la taille de l'échantillon, la marge d'erreur augmente.
    \choice Si on diminue le niveau de confiance, la marge d'erreur augmente.
    \correctchoice Si la variabilité des observations augmente, la marge d'erreur augmente.
    \choice La marge d'erreur ne dépend pas de la taille de l'échantillon.
\end{checkboxes}

\begin{solution}
La marge d'erreur est de la forme $\text{quantile}\times \dfrac{s}{\sqrt{n}}$. Si la variabilité ($s$ ou $\sigma$) augmente, la marge d'erreur augmente. Augmenter $n$ diminue la marge d'erreur. Diminuer le niveau de confiance diminue le quantile, donc la marge d'erreur.
\end{solution}


\question[3] Quelle affirmation est \textbf{vraie} parmi les suivantes?

\begin{checkboxes}
    \choice L'intervalle de confiance pour $\mu$ basé sur la loi de Student est toujours plus étroit que celui basé sur la loi normale centrée réduite.
    \choice On utilise la loi $\chi^2$ pour construire un intervalle de confiance pour la moyenne $\mu$.
    \choice L'intervalle de confiance pour $\sigma^2$ basé sur la loi $\chi^2$ est toujours symétrique autour de $s^2$.
    \correctchoice Lorsque le nombre de degrés de liberté est grand, la loi de Student se rapproche de la loi normale centrée réduite.
\end{checkboxes}

\begin{solution}
La loi $t_\nu$ converge vers $\mathcal{N}(0,1)$ quand $\nu\to\infty$. L'IC basé sur Student est plus large (pas plus étroit) car $t_{\alpha/2;\,n-1}>z_{\alpha/2}$. La loi $\chi^2$ sert pour $\sigma^2$, et l'IC pour $\sigma^2$ n'est pas symétrique car la loi $\chi^2$ est asymétrique.
\end{solution}


\pageblanche
\question\label{varconnue} Soit $X$, le poids (en grammes) de comprimés pharmaceutiques produits par une machine, où $X\sim\mathcal{N}(\mu, \, 4)$. Un échantillon de 25 comprimés donne une moyenne de $\overline{x}=500{,}6$ g.

\begin{parts}
\part[7] Construisez un intervalle de confiance de niveau 95\% pour $\mu$.

\begin{solution}
On a $\sigma^2=4$, donc $\sigma=2$, $n=25$, $\overline{x}=500{,}6$ et $z_{0{,}025}=1{,}96$.

L'IC est:
$$\overline{x}\pm z_{\alpha/2}\frac{\sigma}{\sqrt{n}}=500{,}6\pm 1{,}96\times\frac{2}{\sqrt{25}}=500{,}6\pm 0{,}784=[499{,}816;\, 501{,}384]$$
\end{solution}

\vspace{11cm}

\part[4] Quelle est la taille minimale de l'échantillon à collecter pour que la marge d'erreur ne dépasse pas 0,5 g, en conservant un niveau de confiance de 95\% ?

\begin{solution}
La marge d'erreur est $M=z_{\alpha/2}\dfrac{\sigma}{\sqrt{n}}\le 0{,}5$.

$$n\ge \left(\frac{z_{\alpha/2}\,\sigma}{M}\right)^2=\left(\frac{1{,}96\times 2}{0{,}5}\right)^2=(7{,}84)^2=61{,}47$$

On arrondit vers le haut: $n\ge 62$.
\end{solution}

\end{parts}


\pageblanche
\question\label{varinconnue} Un biologiste mesure la concentration de chlorophylle (en mg/L) dans un lac. Il prélève un échantillon de $n=10$ mesures et obtient:

$$\overline{x}=12{,}3 \quad\text{et}\quad s=1{,}8$$

On suppose que les concentrations suivent une loi normale.

\begin{parts}
\part[7] Construisez un intervalle de confiance de niveau 99\% pour la concentration moyenne $\mu$ de chlorophylle dans ce lac.

\begin{solution}
La variance est inconnue, on utilise la loi de Student. On a $n=10$, donc $\nu=n-1=9$ degrés de liberté. On cherche $t_{0{,}005;\,9}$ dans la table de Student.

$t_{0{,}005;\,9}=3{,}250$ (à vérifier dans la table).

L'IC est:
$$\overline{x}\pm t_{\alpha/2;\,n-1}\frac{s}{\sqrt{n}}=12{,}3\pm 3{,}250\times\frac{1{,}8}{\sqrt{10}}=12{,}3\pm 3{,}250\times 0{,}5692=12{,}3\pm 1{,}850$$
$$=[10{,}450;\, 14{,}150]$$
\end{solution}

\vspace{11cm}

\part[3] Si on avait plutôt utilisé un niveau de confiance de 90\% (au lieu de 99\%), l'intervalle obtenu serait-il plus large ou plus étroit? Justifiez brièvement.

\begin{solution}
L'intervalle serait plus \textbf{étroit}. Un niveau de confiance plus faible signifie un quantile $t_{\alpha/2;\,n-1}$ plus petit, ce qui réduit la marge d'erreur et donc la largeur de l'intervalle.
\end{solution}

\end{parts}


\pageblanche
\question\label{variance} Un ingénieur veut vérifier la stabilité d'une machine-outil. Il mesure le diamètre de 20 pièces produites et obtient une variance échantillonnale de $s^2=0{,}012$ cm$^2$. On suppose que les diamètres suivent une loi normale.

\begin{parts}
\part[8] Construisez un intervalle de confiance de niveau 95\% pour la variance $\sigma^2$.

\begin{solution}
On a $n=20$, donc $\nu=n-1=19$ degrés de liberté, et $s^2=0{,}012$.

On cherche $\chi^2_{0{,}025;\,19}$ et $\chi^2_{0{,}975;\,19}$ dans la table du $\chi^2$.

$\chi^2_{0{,}025;\,19}=32{,}852$ et $\chi^2_{0{,}975;\,19}=8{,}907$.

L'IC pour $\sigma^2$ est:

$$\left[\frac{(n-1)s^2}{\chi^2_{\alpha/2;\,n-1}};\, \frac{(n-1)s^2}{\chi^2_{1-\alpha/2;\,n-1}}\right]=\left[\frac{19\times 0{,}012}{32{,}852};\, \frac{19\times 0{,}012}{8{,}907}\right]=\left[\frac{0{,}228}{32{,}852};\, \frac{0{,}228}{8{,}907}\right]$$
$$=[0{,}00694;\, 0{,}02560]$$
\end{solution}

\vspace{13cm}

\part[3] Déduisez un intervalle de confiance de niveau 95\% pour l'écart-type $\sigma$.

\begin{solution}
On prend la racine carrée des bornes de l'IC pour $\sigma^2$:
$$\text{IC}(\sigma)=\left[\sqrt{0{,}00694};\, \sqrt{0{,}02560}\right]=[0{,}0833;\, 0{,}1600]\;\text{cm}$$
\end{solution}

\end{parts}


\pageblanche
\question\label{proportion} Un sondage effectué auprès de $n=400$ résidents d'une ville révèle que 260 d'entre eux sont en faveur d'un nouveau projet de transport en commun.

\begin{parts}
\part[2] Quelle est l'estimation ponctuelle $\hat{p}$ de la proportion $p$ de résidents favorables au projet?

\begin{solution}
$\hat{p}=\dfrac{260}{400}=0{,}65$.
\end{solution}

\vspace{3cm}

\part[7] Construisez un intervalle de confiance de niveau 95\% pour $p$.

\begin{solution}
On a $\hat{p}=0{,}65$, $n=400$ et $z_{0{,}025}=1{,}96$.

On vérifie que $n$ est assez grand: $n\hat{p}=260\ge 5$ et $n(1-\hat{p})=140\ge 5$. $\checkmark$

L'IC est:
$$\hat{p}\pm z_{\alpha/2}\sqrt{\frac{\hat{p}(1-\hat{p})}{n}}=0{,}65\pm 1{,}96\sqrt{\frac{0{,}65\times 0{,}35}{400}}=0{,}65\pm 1{,}96\times 0{,}02385$$
$$=0{,}65\pm 0{,}04675=[0{,}603;\, 0{,}697]$$
\end{solution}

\vspace{13cm}

\part[3] À partir de cet intervalle, peut-on affirmer, avec 95\% de confiance, que plus de 60\% des résidents sont favorables au projet? Justifiez.

\begin{solution}
Oui. La borne inférieure de l'IC est $0{,}603>0{,}60$. Puisque l'intervalle entier est au-dessus de 0,60, on peut affirmer avec 95\% de confiance que plus de 60\% des résidents sont favorables.
\end{solution}

\end{parts}


\pageblanche
\question Déterminez si les énoncés suivants sont vrais ou faux.

\begin{parts}
\part[2] Lorsqu'on double la taille de l'échantillon, la marge d'erreur d'un intervalle de confiance pour $\mu$ (avec $\sigma$ connu) est divisée par 2.

\begin{oneparcheckboxes}
    \choice Vrai 
    \correctchoice Faux
\end{oneparcheckboxes} \vspace{0.5cm}

\begin{solution}
Faux. La marge d'erreur est proportionnelle à $1/\sqrt{n}$. Si $n$ double, la marge est divisée par $\sqrt{2}\approx 1{,}414$, et non par 2.
\end{solution}

\part[2] Un intervalle de confiance à 99\% est toujours plus large qu'un intervalle de confiance à 95\%, si les mêmes données sont utilisées.

\begin{oneparcheckboxes}
    \correctchoice Vrai 
    \choice Faux
\end{oneparcheckboxes} \vspace{0.5cm}

\begin{solution}
Vrai. Un niveau de confiance plus élevé nécessite un quantile plus grand, ce qui augmente la marge d'erreur et donc la largeur de l'intervalle.
\end{solution}

\part[2] L'intervalle de confiance pour $\sigma^2$ basé sur la loi du $\chi^2$ est symétrique autour de $s^2$.

\begin{oneparcheckboxes}
    \choice Vrai 
    \correctchoice Faux
\end{oneparcheckboxes} \vspace{0.5cm}

\begin{solution}
Faux. La loi $\chi^2$ est asymétrique, donc l'IC pour $\sigma^2$ n'est pas symétrique autour de $s^2$.
\end{solution}

\part[2] Pour construire un intervalle de confiance pour la proportion $p$, on a besoin que la taille d'échantillon soit suffisamment grande.

\begin{oneparcheckboxes}
    \correctchoice Vrai 
    \choice Faux
\end{oneparcheckboxes} \vspace{0.5cm}

\begin{solution}
Vrai. L'IC pour $p$ repose sur l'approximation normale de la loi binomiale, qui nécessite $n\hat{p}\ge 5$ et $n(1-\hat{p})\ge 5$.
\end{solution}

\part[2] Si $X_1,\ldots,X_n$ proviennent d'une loi inconnue et que $n$ est grand, on peut construire un intervalle de confiance \textbf{approximatif} pour $\mu$ en utilisant le théorème central limite.

\begin{oneparcheckboxes}
    \correctchoice Vrai 
    \choice Faux
\end{oneparcheckboxes} \vspace{0.5cm}

\begin{solution}
Vrai. Le TCL garantit que $\overline{X}\approx\mathcal{N}(\mu,\sigma^2/n)$ pour $n$ grand, ce qui permet de construire un IC approximatif.
\end{solution}

\part[2] La longueur d'un intervalle de confiance symétrique pour $\mu$ est égale à deux fois la marge d'erreur.

\begin{oneparcheckboxes}
    \correctchoice Vrai 
    \choice Faux
\end{oneparcheckboxes} \vspace{0.5cm}

\begin{solution}
Vrai. Un IC symétrique est de la forme $\overline{x}\pm M$, donc sa longueur est $(\overline{x}+M)-(\overline{x}-M)=2M$.
\end{solution}

\end{parts}


\pageblanche
\question\label{TCL} Le temps d'attente (en minutes) au guichet d'une banque suit une loi de moyenne $\mu=7$ et d'écart-type $\sigma=3$ (la forme de la distribution est inconnue). On observe un échantillon de $n=36$ clients.

\begin{parts}
\part[3] Justifiez pourquoi on peut construire un intervalle de confiance pour $\mu$ malgré le fait que la loi de la population n'est pas normale.

\begin{solution}
Puisque $n=36\ge 30$, on peut invoquer le théorème central limite (TCL): $\overline{X}\approx\mathcal{N}(\mu,\sigma^2/n)$. Cela permet de construire un IC approximatif.
\end{solution}

\vspace{5cm}

\part[7] L'échantillon donne $\overline{x}=7{,}8$ et $s=3{,}2$. Construisez un intervalle de confiance de niveau 90\% pour $\mu$.

\begin{solution}
Par le TCL, on utilise l'IC approximatif avec $z_{\alpha/2}=z_{0{,}05}=1{,}645$.

$$\overline{x}\pm z_{\alpha/2}\frac{s}{\sqrt{n}}=7{,}8\pm 1{,}645\times\frac{3{,}2}{\sqrt{36}}=7{,}8\pm 1{,}645\times 0{,}5333=7{,}8\pm 0{,}877$$
$$=[6{,}923;\, 8{,}677]$$
\end{solution}

\end{parts}


\pageblanche

\question[5] Un échantillon de taille $n=12$ provenant d'une population normale donne $\overline{x}=45$ et $s=6$. On calcule un intervalle de confiance pour $\mu$ à un certain niveau de confiance et on obtient $[45-3{,}81;\, 45+3{,}81]=[41{,}19;\, 48{,}81]$. Quel est le niveau de confiance de cet intervalle? Montrez votre démarche.

\begin{solution}
La variance est inconnue et $n=12$, donc on utilise Student avec $\nu=11$ degrés de liberté.

La marge d'erreur est $M=3{,}81=t_{\alpha/2;\,11}\times\dfrac{6}{\sqrt{12}}=t_{\alpha/2;\,11}\times 1{,}7321$.

Donc $t_{\alpha/2;\,11}=\dfrac{3{,}81}{1{,}7321}\approx 2{,}201$.

En consultant la table de Student à 11 degrés de liberté, $t_{0{,}025;\,11}=2{,}201$.

Donc $\alpha/2=0{,}025$, soit $\alpha=0{,}05$, et le niveau de confiance est $1-\alpha=0{,}95$ (95\%).
\end{solution}


\end{questions}
\end{document}
